\scanpage{69}
c) \deleted{Soient} $\vec D,\vec D',\vec\Delta$ ps. droites \inserted{orientées à supp. distincts}, $x,x'$ les
int. de $\Delta$ avec $D,D'$ (on ne suppose pas $x\ne x'$)
et $s,s'$ des sommets sur $\Delta$ distincts de $x,x'$.
Supposons $\rho_{\vec D\unsure{\,|,}s}=\rho_{\vec D',s}$. Pour qu'on ait
$\rho_{\vec D,s'}=\rho_{\vec D',s'}$, il f. et s. que $s,s'$ \unsure{appart\textsuperscript{nt}|appartiennent}
au \unsure{m\textsuperscript{e}|$=$} arc de $\Delta$ déterminé par $x,x'$

\textcircled{d} Soit $\vec\Delta_1\in\unsure{\vec{\mathscr D}|\vec{\mathcal D}}(s)$, $\vec D\in\vec{\mathscr D}\setminus\vec{\mathscr D}(s)$, $\vec\Delta^{\deleted{\unsure{1}}}_2=\rho_{\vec D,s}(\vec\Delta)$,
Soient $\{u_1\}=|\vec\Delta_1|\cap|\vec D|$, $\{u^{\unsure{\prime|{}}}_2\}\unsure{=|\Leftarrow}|\vec\Delta_2|\cap|\vec D|$, alors $u'_2$ est le
sommet sur $\vec D$ qui suit $u_1$, \uuline{parmi} toutes les
intersections de $\vec D$ avec les ps. dr. passant par $s$.

e)
% archive sheet number: 68
